% ---------------------------------------------------------------
% Section 6: Integration with the BX3 Framework
% ---------------------------------------------------------------
\section{Integration with the BX3 Framework}
\label{sec:bx3-integration}

The Recursive Spawning protocol gains its formal structure and
correctness guarantees through deep integration with the BX3 Framework.
BX3 supplies the three-tier epistemic architecture --- Purpose Layer,
Fact Layer, and Projection Layer --- that anchors spawning decisions,
enforces depth bounds, and maintains a forensic ledger of every
spawning event. This section details how each layer participates in
the protocol and how the Bounds Engine operates within this framework.

\subsection{Purpose Layer as Accountability Anchor}
\label{sec:purpose-layer-anchor}

The \textbf{Purpose Layer} is the single source of truth from which all
spawned agents derive their legitimate scope of action. When a spawning
event occurs, the parent agent records the originating purpose
identifier (PID) alongside the spawn request. The child agent inherits this
PID as a non-transferable attribute encoded in its header and verifiable
against the Purpose Layer at any point during execution.

This inheritance model has three consequences for spawning behavior:

\begin{enumerate}
  \item \textbf{Scope monotonicity:} The child's admissible action set is
        a subset of the parent's admissible set, as constrained by the
        PID hierarchy. No spawned agent may expand scope beyond the
        purpose that justified its creation.
  \item \textbf{Accountability traceability:} Every action taken by a
        spawned agent traces back through its chain of parent PIDs to a
        single root purpose at the time of top-level initialization. The
        Purpose Layer provides an O(1) lookup for the chain root given any
        descendant PID.
  \item \textbf{Cascade termination anchor:} The Purpose Layer enforces
        a special \texttt{TERMINATE} status for PIDs whose associated
        purpose has been fulfilled, invalidated, or revoked. When a PID
        enters \texttt{TERMINATE} state, no further spawning is permitted
        from any agent holding that PID or any of its descendants. This
        prevents orphaned cascades from continuing beyond their
        productive horizon.
\end{enumerate}

In the Recursive Spawning protocol, the Purpose Layer is queried at
spawn initiation (\S\ref{sec:spawn-initiation}) and at every subsequent
depth-bounded iteration via the \texttt{check\_purpose\_validity(p)}
primitive. This primitive returns \texttt{true} only if the PID is in
an \texttt{ACTIVE} state and its depth constraint has not been violated.

\subsection{Bounds Engine and Depth Enforcement}
\label{sec:bounds-engine}

The \textbf{Bounds Engine} is the run-time component responsible for
enforcing the spawning depth constraint $D_{\max}$ defined in the
protocol parameters. It operates as a layer-local service within the
BX3 Framework and maintains the following state:

\begin{itemize}
  \item $\texttt{current\_depth}(p)$: the number of generations from
        the root PID to agent $p$.
  \item $\texttt{budget}(p)$: the remaining spawn budget available to
        agent $p$, initialized to $B$ at each root spawn and decremented
        at each successful child creation.
  \item $\texttt{path\_signature}(p)$: an immutable cryptographic hash
        chain linking $p$ to its chain root, used to detect depth frauds.
\end{itemize}

The Bounds Engine exposes two primitives to the spawning protocol:

\begin{enumerate}
  \item $\texttt{check\_depth\_validity}(p)$: returns \texttt{true} if
        $\texttt{current\_depth}(p) < D_{\max}$ and $\texttt{budget}(p) > 0$.
  \item $\texttt{record\_spawn}(p_{\text{parent}}, p_{\text{child}})$:
        atomically increments the depth of the child, decrements the
        parent's budget, and extends the path signature chain with the
        child's PID.
\end{enumerate}

Both primitives are implemented as atomic transactions against the BX3
Framework's shared state store. The depth validity check is performed
before any spawning fork is committed, ensuring that a spawn request
with $\texttt{current\_depth}(p_{\text{parent}}) \geq D_{\max}$ is
rejected at the Purpose Layer without requiring a full forensic ledger
write. This early-exit optimization prevents wasted computational
resources on malformed spawns.

The Bounds Engine also enforces the branching factor constraint $B$,
which bounds the number of direct children any agent may spawn. When
$\texttt{budget}(p) = 0$, the agent may not spawn additional children
regardless of remaining depth capacity. This dual-constraint model
($D_{\max}$ and $B$) ensures that the total number of agents in any
chain is bounded by $O(B^{D_{\max}})$, enabling finite-space reasoning
about termination and safety properties.

\subsection{Fact Layer as Forensic Ledger}
\label{sec:fact-layer}

The \textbf{Fact Layer} serves as the immutable, append-only forensic
ledger of all spawning events in the BX3 Framework. Every spawn
initiation, depth rejection, purpose revocation, or cascade termination
event is recorded in the Fact Layer as a \texttt{spawn\_record} with the
following schema:

\begin{verbatim}
spawn_record {
    event_id     : UUIDv7
    timestamp    : ISO-8601(μs)
    parent_pid   : PID
    child_pid    : PID | null
    event_type   : INITIATED | COMMITTED | REJECTED_DEPTH
                   | REJECTED_PURPOSE | TERMINATED
    depth        : nat
    purpose_hash : H(purpose_context)
    parent_ref   : H(parent_record)  // immutable parent pointer
}
\end{verbatim}

The \texttt{parent\_ref} field encodes the hash pointer to the parent's
spawn record, forming an immutable chain that is cryptographically
self-certifying. This design ensures that no agent can fabricate a
spawn history: any modification to a historical record breaks the hash
chain, detectable by any auditor with read access to the Fact Layer.

The forensic ledger serves three critical roles in the Recursive
Spawning protocol:

\begin{enumerate}
  \item \textbf{Drift detection:} By comparing the on-chain parent
        pointer of an agent against the expected parent derived from its
        current spawning context, the protocol detects and blocks
        spawned agents that have been redirected or ``drifted'' from
        their original chain of custody.
  \item \textbf{Integrity verification:} The Fact Layer provides a
        verifiable sequence of all spawning events, enabling post-hoc
        auditing of any agent's origin, purpose, and chain of ancestors.
  \item \textbf{Non-repudiation:} Because spawn records are
        append-only and hash-chained, no agent can deny its participation
        in a spawning event after the record has been committed.
\end{enumerate}

The Fact Layer is queried during the initialization phase of any agent
that requires verification of its chain integrity, particularly when
an agent re-enters the system after a suspension or a distributed
checkpoint/restoration event. The integrity check reconstructs the full
ancestry chain from the Fact Layer and verifies that each successive
parent pointer is consistent with the recorded history.

\subsection{Combined Protocol Flow}
\label{sec:combined-flow}

When the three layers operate in concert, the Recursive Spawning
protocol follows the following sequence for each spawn event:

\begin{enumerate}
  \item The parent agent calls \texttt{spawn(child\_purpose)} against
        the BX3 Framework interface.
  \item The Purpose Layer resolves the parent's PID and verifies it is in
        \texttt{ACTIVE} state.
  \item The Bounds Engine evaluates $\texttt{check\_depth\_validity}(p)$.
        If depth or budget constraints are violated, a
        \texttt{REJECTED\_DEPTH} record is written to the Fact Layer and
        the spawn is aborted.
  \item If depth validity holds, the child PID is allocated and the
        \texttt{record\_spawn} primitive extends the path signature chain.
  \item A \texttt{COMMITTED} spawn record is written to the Fact Layer,
        linking the child's PID to the parent's via an immutable hash
        pointer.
  \item The child agent initializes with the inherited PID, depth
        incremented by one, and budget set to $B$.
\end{enumerate}

This sequence ensures that every spawning event is simultaneously
authorized by the Purpose Layer, constrained by the Bounds Engine, and
recorded immutably in the Fact Layer. The three-layer architecture thus
provides a complete formal foundation for the Recursive Spawning
protocol's correctness guarantees.

% ---------------------------------------------------------------
% Section 7: Correctness Properties
% ---------------------------------------------------------------
\section{Correctness Properties}
\label{sec:correctness}

This section establishes the formal correctness guarantees provided by
the Recursive Spawning protocol under the BX3 Framework. We prove three
central properties: (1) the prevention of agent drift, (2) the integrity
of the spawn chain, and (3) the guaranteed termination of all spawn
chains within bounded depth. Together, these properties ensure that
the protocol behaves as a safe, predictable, and auditable mechanism
for recursive agent creation.

\subsection{Drift Prevention}
\label{sec:drift-prevention}

\begin{definition}[Drift]
An agent $a$ is said to have \textbf{drifted} if its current parent
pointer $p_a$ does not equal the authentic parent derived from the
immutable forensic ledger of spawn records $\mathcal{F}$.
\end{definition}

Drift can occur when a malicious or faulty intermediate agent redirects
a spawned child to a different parent than the one recorded at spawn
commitment. The Recursive Spawning protocol prevents drift through the
interaction of three mechanisms:

\begin{enumerate}
  \item \textbf{Immutable parent pointer:} At spawn commitment, the
        child's record in the Fact Layer is written with the hash of the
        parent's record as \texttt{parent\_ref}. This pointer is
        cryptographically immutable: modifying the parent's record breaks
        the hash chain, making the drift immediately detectable.
  \item \textbf{Forensic ledger verification:} During any agent
        initialization or re-entry, the protocol queries the Fact Layer
        for the agent's \texttt{parent\_ref} and compares it against the
        claimed parent's record hash. A mismatch triggers a
        \texttt{DRIFT\_DETECTED} exception and halts the agent's
        execution.
  \item \textbf{Purpose Layer chain termination:} The Purpose Layer
        enforces that every agent's PID is linked to a chain root in an
        \texttt{ACTIVE} state. An agent whose parent chain has been
        broken cannot re-validate its PID against the Purpose Layer,
        and therefore cannot continue execution.
\end{enumerate}

\begin{Theorem}[Drift Prevention]
No agent spawned under the Recursive Spawning protocol can exhibit
drift.
\end{Theorem}

\begin{Proof}
Consider an agent $a$ spawned with parent $p$. At spawn commitment,
the Fact Layer records \texttt{parent\_ref}($a$) = $H(\text{record}(p))$.
By construction of the hash chain, $H(\text{record}(p))$ is a function of
the complete contents of $p$'s spawn record, including all prior
parent pointers in its own chain. Any modification to $p$'s record
changes $H(\text{record}(p))$ and invalidates the chain.

Assume, for contradiction, that $a$ drifts. Then there exists an
intermediate agent $m$ that, after the spawn of $a$, changes the
recorded parent of $a$ to a different agent $p' \neq p$. For this to
occur, $m$ must modify the Fact Layer record of $a$, which is append-only;
modification would break the hash chain of any agent that subsequently
reads $a$'s record. The forensic ledger's immutability property is
guaranteed by the BX3 Framework's append-only store, which rejects
in-place updates to committed records. Therefore, no drift can be
committed without detection.

Furthermore, the Purpose Layer authorization performed at every spawn
event validates that the parent PID is in an \texttt{ACTIVE} state before
committing the spawn record. This prevents drift via ``ghost parents''
(PIDs that do not correspond to live agents). ∎
\end{Proof}

\subsection{Chain Integrity}
\label{sec:chain-integrity}

\begin{definition}[Spawn Chain Integrity]
A spawn chain $C = (a_0, a_1, \dots, a_n)$ satisfies \textbf{chain
integrity} if for every $i \in \{1, \dots, n\}$, the parent pointer of
$a_i$ in the Fact Layer resolves to $a_{i-1}$.
\end{definition}

Chain integrity ensures that ancestry records in the Fact Layer correctly
reflect the true genealogical structure of the spawn graph, without
gaps, loops, or broken links.

\begin{Theorem}[Chain Integrity Invariant]
At all times during protocol execution, the spawn chain of any agent
satisfies chain integrity.
\end{Theorem}

\begin{Proof}
The invariant holds vacuously at initialization: the root agent $a_0$
has no parent, and its chain consists of a single element.

Assume the invariant holds for a chain of length $k$. Consider the
spawning of $a_{k+1}$ from $a_k$. The spawn protocol requires:

\begin{enumerate}
  \item Purpose Layer validation that $\text{PID}(a_k)$ is \texttt{ACTIVE}.
  \item Bounds Engine confirmation that $\text{depth}(a_k) < D_{\max}$.
  \item Atomic execution of $\texttt{record\_spawn}(a_k, a_{k+1})$,
        which writes \texttt{parent\_ref}($a_{k+1}$) = $H(\text{record}(a_k))$.
\end{enumerate}

Step 3 atomically establishes the parent link in the Fact Layer.
Because the write is atomic and the Fact Layer is append-only, the
parent pointer of $a_{k+1}$ is guaranteed to resolve to $a_k$ in the
ledger. Since the prior chain of length $k$ satisfies the invariant by
the inductive hypothesis, the extended chain of length $k+1$ also
satisfies it. ∎

The chain integrity invariant is preserved under agent suspension and
restoration: upon re-entry, an agent's ancestry is reconstructed from
the Fact Layer and verified against the chain integrity invariant before
execution resumes. If the invariant is violated (e.g., due to a node
failure during an atomic write), the agent enters a \texttt{INVALID\_CHAIN}
state and awaits manual resolution.

\subsection{Termination}
\label{sec:termination}

\begin{Theorem}[Bounded Termination]
Every spawn chain initiated under the Recursive Spawning protocol
terminates within at most $D_{\max}$ generations.
\end{Theorem}

\begin{Proof}
Each agent $a$ at depth $d$ satisfies $d < D_{\max}$ to be permitted
to spawn. By the $\texttt{check\_depth\_validity}$ primitive, an agent
at depth $d = D_{\max} - 1$ may spawn one final generation of children,
each at depth $D_{\max}$. Agents at depth $D_{\max}$ are subject to the
following constraints:

\begin{enumerate}
  \item $\texttt{check\_depth\_validity}(a)$ returns \texttt{false},
        preventing further spawning by the Bounds Engine.
  \item The Purpose Layer enforces that the PID of any agent at depth
        $D_{\max}$ remains in \texttt{ACTIVE} state (the Purpose Layer
        does not terminate agents based on depth alone), but no further
        spawn operations are actionable from depth $D_{\max}$.
  \item The branching factor bound $B$ ensures that even at depth
        $D_{\max} - 1$, at most $B$ children are generated, capping the
        total number of agents in any chain at $B^{D_{\max}}$.
\end{enumerate}

Therefore, after at most $D_{\max}$ generations, no agent in the chain
can invoke the spawn primitive successfully. The chain reaches a
natural termination state where all agents have exhausted both their
depth budget and their branching budget. ∎
\end{Proof}

The termination theorem guarantees that the Recursive Spawning protocol
does not admit infinite descent. This property is critical for any
practical deployment: without a depth bound, the protocol would be
capable of generating unbounded numbers of agents in finite time, which
would violate resource-bounded execution guarantees required in any
real compute environment.

Together, drift prevention, chain integrity, and bounded termination
constitute the core correctness guarantees of the Recursive Spawning
protocol under the BX3 Framework. These properties ensure that spawning
behaves as a safe, accountable, and predictable primitive suitable for
use in production agentic systems.